\chapter{Sustentabilidade fiscal: regras, instituições e evidências}
\label{chap:revisao}

\section{Regras fiscais, desenho institucional e aplicação subnacional}
\label{sec:regras}

As regras fiscais ocupam papel central na literatura de finanças públicas por buscarem limitar o viés deficitário, a fim de conter o acúmulo excessivo de dívida e fortalecer a credibilidade do governo. Em diferentes contextos, a administração pública pode enfrentar incentivos para elevar gastos, reduzir tributos ou adiar ajustes, principalmente quando os benefícios políticos dessas decisões são imediatos e seus custos fiscais só se farão sentir no futuro. Esse problema é formalizado por \citeonline{alesina1990}, que mostram como a alternância de governos com preferências distintas pode levar ao uso estratégico da dívida pública, gerando déficits elevados.

A fim de conter a discricionariedade da política fiscal, a literatura destaca o uso das regras fiscais, tipicamente divididas pelo agregado que buscam controlar. Cada tipo atua sobre uma margem distinta: regras de dívida limitam o estoque acumulado, mas não impedem déficits enquanto o teto não é atingido; regras de receita estabelecem tetos para a carga tributária ou vinculam excessos de arrecadação, embora fiquem vulneráveis a choques de crescimento; regras de resultado atuam sobre o fluxo anual, embora sejam mais sensíveis ao ciclo econômico; regras de despesa restringem o crescimento do gasto, podendo penalizar o investimento quando o ajuste recai sobre rubricas discricionárias. No plano subnacional, ganham peso regras de operações de crédito e serviço de dívida, essenciais para condicionar o endividamento dos entes.

A efetividade dessas regras, entretanto, depende da capacidade de coordenação entre níveis de governo, que é estruturalmente mais complexa no caso subnacional. Estados e governos locais possuem políticas próprias de gasto, condicionadas em grande parte por transferências intergovernamentais, autorizações legais para contratação de crédito e limites impostos pela autoridade central. Dada essa interdependência, regras fiscais subnacionais servem não apenas para disciplinar a política fiscal local, mas também para preservar a estabilidade fiscal da federação como um todo. Como destacam \citeonline{liu2011}, marcos de responsabilidade fiscal podem ajudar a coordenar diferentes níveis de governo e mitigar comportamentos fiscalmente irresponsáveis que geram benefícios de curto prazo, mas custos coletivos no médio e longo prazo.

Em federações, essa discussão se conecta a problemas de restrição orçamentária branda \cite{kornaiSoftBudgetConstraint1986} e de fragmentação fiscal \cite{velasco2000}. \citeonline{rodden2002} argumenta que a expectativa de apoio do governo central em momentos de crise pode reduzir os incentivos ao ajuste fiscal preventivo por parte de governos subnacionais. Esse problema também se relaciona ao uso comum de recursos fiscais, uma vez que os benefícios políticos do gasto podem ser apropriados localmente, enquanto parte dos custos é transferida para o restante da federação. 

Regras fiscais, porém, não funcionam automaticamente. Como discutem \citeonline{wyplosz2012}, \citeonline{heinemann2018}, \citeonline{caselli2020} e \citeonline{brandleFiscalRulesMatter2024}, seus efeitos dependem do desenho institucional, da abrangência dos agregados controlados, da transparência das contas públicas, da existência de mecanismos de enforcement, da credibilidade de sua aplicação e de sanções efetivas para o descumprimento. Regras mal calibradas, pouco críveis ou sujeitas a sucessivas exceções podem ter efeitos limitados sobre o comportamento fiscal dos governos.

A evidência internacional confirma essa interpretação. \citeonline{biase2022} e \citeonline{vammalle2021} documentam que regras fiscais subnacionais se expandiram em países da OCDE, mas que regras formalmente semelhantes produzem graus de rigidez efetiva distintos conforme cobertura, monitoramento e mecanismos de exceção. Para os estados norte-americanos, \citeonline{poterba1994} mostra que instituições orçamentárias influenciam a velocidade de ajuste diante de crises fiscais. \citeonline{venturiniUnintendedCompositionEffect2020}, analisando municípios italianos, encontra que regras mais restritivas podem recompor a despesa contra o investimento público — padrão documentado também para países latino-americanos por \citeonline{calderon2003} —, sugerindo um efeito não intencional potencialmente relevante para o contexto subnacional brasileiro.

No caso dos estados brasileiros, os contratos de refinanciamento da Lei n\textsuperscript{o}~9.496/1997 aproximam-se de uma regra de serviço da dívida, condicionando o fluxo anual de pagamentos dos estados à União. A diversidade dos limites contratuais entre os estados é explorada no Capítulo~\ref{chap:institucional}.


\section{Restrição orçamentária intertemporal e função de reação fiscal}

A discussão sobre regras fiscais conduz diretamente ao problema que essas regras buscam endereçar: a sustentabilidade da dívida pública. Se, de um lado, tais regras tentam mudar os incentivos dos participantes, buscando induzir comportamentos fiscais diferentes, de outro, sua efetividade é avaliada, em última análise, pela capacidade de contribuir para trajetórias fiscais compatíveis com a solvência intertemporal.

A sustentabilidade fiscal está associada à capacidade do governo de honrar seus compromissos ao longo do tempo sem recorrer ilimitadamente à emissão de nova dívida para financiar obrigações passadas. Apesar de haver várias definições, em termos gerais, uma trajetória fiscal é considerada sustentável quando políticas de gasto, arrecadação e endividamento são compatíveis com a restrição orçamentária intertemporal (ROI) do governo, ou seja, exige-se que o endividamento não cresça de forma explosiva em relação à capacidade de pagamento do setor público.

A literatura analisou essa questão a partir da ROI, indicando a necessidade do valor presente dos superávits primários futuros ser suficiente para cobrir o estoque inicial da dívida. Antes do desenvolvimento dos testes econométricos de sustentabilidade, a avaliação da solvência era frequentemente conduzida por indicadores contábeis e pela análise da dinâmica da dívida, instrumentos úteis para descrever pressões fiscais correntes, mas com pouca capacidade de verificação no sentido intertemporal.

A partir dessa limitação, uma primeira geração de estudos passou a avaliar a solvência fiscal por meio das propriedades estatísticas das séries fiscais. Trabalhos como \citeonline{hamilton1985}, \citeonline{trehan1988} e \citeonline{hakkio1991} utilizaram testes de estacionariedade e cointegração para investigar se receitas, despesas e dívida pública apresentavam comportamento compatível com a ROI. Esses testes tornaram a sustentabilidade fiscal empiricamente testável, mas apresentavam uma limitação importante: não revelam como o governo reage ao aumento do endividamento. Em particular, não indicam se a autoridade fiscal responde ao crescimento da dívida por meio de maior esforço primário.

Essa distinção é relevante porque a trajetória observada da dívida pode refletir fatores fora do controle fiscal corrente, como variações nas taxas de juros, crescimento econômico, choques na receita e alterações institucionais em geral. No caso de governos subnacionais, tal limitação é ainda mais evidente, visto que a dívida estadual pode ser afetada por variáveis que dependem do governo central. Assim, a análise da sustentabilidade da dívida estadual exige observar não apenas o comportamento da dívida, mas também se há uma resposta fiscal sistemática ao seu aumento.

A fim de lidar com essas limitações, \citeonline{bohn1998} propõe uma abordagem alternativa, baseada na função de reação fiscal, investigando a reação do resultado primário frente ao aumento do endividamento em vez de avaliar apenas as propriedades estatísticas da dívida. A intuição é que, se o governo eleva sistematicamente seu esforço fiscal quando a dívida cresce, há evidência compatível com a existência de mecanismos fiscais de estabilização da dívida. Formalmente, a função de reação fiscal assume a forma:

\begin{equation}
pb_{t} = \alpha + \rho \, d_{t-1} + \mathbf{X}_t\boldsymbol{\gamma} + \varepsilon_t
\label{eq:bohn}
\end{equation}

\noindent onde $pb_{t}$ é o resultado primário, $d_{t-1}$ é a dívida defasada e $\mathbf{X}_t$ é um vetor de controles — em \citeonline{bohn1998}, representado por medidas de gastos governamentais temporários ($GVAR_t$) e de ciclo econômico ($YVAR_t$) para o contexto americano. A condição de sustentabilidade requer $\rho > 0$: o resultado primário responde positivamente ao aumento da dívida, sinalizando um mecanismo de correção fiscal. No presente trabalho, essa especificação é estendida ao contexto de painel, com a adição de efeitos fixos de estado e ano e controles adaptados ao contexto subnacional brasileiro, conforme detalhado no Capítulo~\ref{chap:metodologia}.

A abordagem foi posteriormente reforçada por \citeonline{bohn2007}, que argumenta que restrições de estacionariedade e cointegração não constituem condições necessárias para que a política fiscal seja compatível com a restrição orçamentária intertemporal. O resultado é formalmente robusto, mas sua aplicação empírica exige cautela: as conclusões são condicionais à especificação adotada, à amostra analisada e aos controles utilizados, podendo a magnitude da resposta fiscal variar conforme o ambiente institucional em que a política fiscal é conduzida.

No caso brasileiro, \citeonline{mello2005} aplica essa abordagem ao setor público e a diferentes níveis de governo, encontrando evidência de reação positiva ao aumento do endividamento, com fortalecimento dessa resposta após o período de renegociação e ajuste no final dos anos 1990. \citeonline{luporiniSustainabilityBrazilianFiscal2015}, com coeficientes variantes no tempo, encontrou que a política fiscal brasileira se tornou mais estável após 2000, mas progressivamente menos responsiva ao nível da dívida. 

Mais recentemente, a literatura passou a incorporar variáveis institucionais diretamente aos testes de sustentabilidade. Em vez de avaliar apenas se o resultado primário reage à dívida, esses estudos investigam se regras fiscais, programas de alívio da dívida ou características institucionais condicionam tal reação. Paralelamente, avançou-se em direções complementares, incluindo testes de fadiga fiscal e limites de dívida, modelos dinâmicos de painel para a trajetória do endividamento e especificações que permitem respostas fiscais diferenciadas conforme o nível de dívida. 


\section{Evidências empíricas}

A literatura empírica sobre regras fiscais e sustentabilidade da dívida apresenta resultados que variam entre países, níveis de governo e arranjos institucionais. Em alguns contextos, regras fiscais e instituições orçamentárias contribuem para respostas compatíveis com a estabilização da dívida, ao passo que em outros, a existência formal de regras não é suficiente para induzir uma reação fiscal adequada, sobretudo quando sua aplicação é pouco crível ou quando há expectativa recorrente de renegociação.

Para os estados norte-americanos, como discutido na Seção~\ref{sec:regras}, as instituições fiscais parecem exercer papel disciplinador. Na mesma linha, \citeonline{mahdaviBohnsTestFiscal2014} aplica o teste de sustentabilidade de Bohn aos governos estaduais norte-americanos e encontra evidência compatível com sustentabilidade fiscal subnacional, sugerindo que, nesse contexto, os estados reagem ao aumento da dívida por meio de ajustes fiscais, indicando mecanismos de estabilização do endividamento.

Esses resultados, porém, não são universais. \citeonline{abubakarFiscalStanceFiscal2025} analisam a relação entre postura fiscal, regras fiscais e dinâmica da dívida pública nos Estados Unidos e mostram que a presença de regras fiscais não foi suficiente para conter a trajetória de crescimento da dívida do governo federal. \citeonline{abubakarDebtReliefFiscal2025}, ao analisar 37 países da África Subsaariana, encontram que programas de alívio da dívida melhoram a sustentabilidade fiscal, mas que a imposição de regras aparece associada à sua piora. 

Evidências de outras federações também apontam para a importância do contexto institucional. Na Índia, \citeonline{renjithDynamicsPublicDebt2020} estimam funções de reação fiscal para 22 estados e encontram que a sustentabilidade se verifica em apenas metade dos entes, evidenciando que regras uniformes produzem resultados assimétricos conforme as condições estruturais de cada estado. No México, \citeonline{cabralSustainabilitySubnationalPublic2022} mostram que a avaliação da solvência subnacional exige atenção às diferenças fiscais e institucionais entre unidades federativas, e \citeonline{delcastilloSubnationalPublicDebt2024} encontram que a nova regra fiscal mexicana contribuiu para reduzir o endividamento subnacional.

No Brasil, os resultados são mistos. \citeonline{bastos2013}, ao estimar funções de reação fiscal e limites de dívida para os estados brasileiros, mostram que o espaço fiscal varia de forma relevante entre unidades federativas. \citeonline{caldeira2016} encontram evidência compatível com sustentabilidade,  identificando relação de longo prazo entre dívida líquida e resultado primário. Por outro lado, \citeonline{tabosa2016} encontram que, em média, os estados brasileiros não apresentavam uma política fiscal suficientemente ativa de geração de superávits diante do aumento da dívida pública.

Estudos mais recentes mantêm esse debate aberto. \citeonline{arraes2024} indicam que a LRF influenciou a gestão fiscal estadual sem eliminar seus desafios estruturais. \citeonline{ribeiroSustentabilidadeFiscalDos2024} constatam uma reação empírica negativa ao acúmulo de passivos, com o saldo primário gerado não tendo sido suficiente para conter a trajetória da dívida, enquanto \citeonline{luduviceSustainabilityStateGovernment2025} encontram evidência de sustentabilidade conjunta da dívida financeira estadual entre 2007 e 2022, ainda que próxima à fronteira entre estabilidade e instabilidade.

Apesar dos avanços dessa literatura, no contexto subnacional brasileiro a maioria dos estudos avalia a existência de sustentabilidade ou o efeito médio de regras fiscais, sem explorar a heterogeneidade institucional dos contratos de refinanciamento. O caso brasileiro oferece uma oportunidade nessa direção, na medida em que os contratos da Lei nº 9.496/1997 estabeleceram condições diferenciadas entre estados — como taxas de juros e limites de comprometimento da receita —, o que o presente trabalho explora como fonte de variação institucional na estimação da função de reação fiscal.

Em conjunto, essas evidências reforçam que a relação entre regras fiscais e sustentabilidade da dívida não é mecânica. Regras fiscais podem contribuir para conter o endividamento quando são críveis, bem desenhadas e acompanhadas de mecanismos efetivos de monitoramento e execução. Contudo, a adoção formal de regras pode ter efeitos limitados, ou mesmo contraproducentes, quando não altera os incentivos fiscais dos governos ou quando está inserida em ambientes de baixa credibilidade.